\section{CoM制御の役割と本研究における位置づけ}

歩行制御におけるCoM制御は，歩行生成と安定化制御をつなぐ中核的な役割を担っている．上位層で生成されたCoM参照軌道に基づき，実際のCoM運動を調整することで，ZMPを支持多角形内に保ち，歩行の安定性を確保する．

従来の多くの手法では，CoMの水平方向の運動制御に主眼が置かれており，鉛直方向の運動は一定もしくは単純な関数として扱われてきた．しかし，床が変形する環境では，CoM高さの変動がZMPや床反力分布に影響を及ぼし，歩行安定性を低下させる要因となる．

本研究では，この点に着目し，床沈み込みに応じてCoM高さ参照を調整することで，歩行安定性を向上させる手法を検討する．本章で述べた基礎概念を踏まえ，次章では既存研究における歩行制御手法とその課題について整理する．